\section{What the hardware takes over}
\label{sec:handover}

CAN solves the same problems as \kapref{ch:frames}--\kapref{ch:bussar}, in silicon:

\begin{booktable}{@{}lL@{}}
\thead{Mechanism} & \thead{Who does it in CAN} \\ \midrule
Frame synchronisation & The controller: the SOF bit, and seven recessive bits as an end
  marker. \\
Length field & The controller: the DLC, four bits, 0--8 data bytes. \\
Checksum & The controller: CRC-15, computed and checked by the same engine. \\
Collision handling & The bus itself, through arbitration --- without a bus master and without
  lost time. \\
Acknowledgement & The controller: a single bit inside the frame, which every receiver pulls
  low. \\
Addressing & \textbf{Does not exist.} The identifier describes the content, and every node
  filters for itself. \\
Sequence number & \textbf{Does not exist.} If one is needed, the application has to put it in the
  payload. \\
Timeout and retransmission & A complete controller does it in hardware. \emph{This one does
  not} --- see \secref{sec:limits}. \\
\end{booktable}

What is left for the software is therefore to hand over an identifier, a length and a few data
bytes --- and to read off how it went.

\begin{aside}
\note CAN as a protocol --- the bus, the frame, bit stuffing, CRC-15, arbitration, bit timing and
error handling --- is covered in full in the course's CAN book. This chapter assumes its chapters
1, 2 and 7. Read those first; the rest of that book is further reading you can take when you need
it.
\end{aside}

\section{The controller, and why a layer sits in between}
\label{sec:layer}

The controller is designed in VHDL by a parallel class. It is built for a caller in the same clock
domain, and therefore signals with pulses that are \textbf{one clock cycle} long --- 20~ns at
50~MHz.

A processor polling over a serial bus sees the system a few tens of microseconds at a time at
best. A 20~ns pulse is therefore gone thousands of cycles before the next poll arrives.

That is why there are two layers between the controller and your code:

\begin{console}
                your C++ code
                      |
                     SPI                  5-byte transactions, <= 1 MHz
                      |
   spi_slave  ->  spi_reg_bridge          bytes become transactions
                      |
              register_bank               pulses become sticky, pollable bits
                      |
              can_controller              CAN in VHDL
                      |
                  the CAN bus
\end{console}

\textbf{The register bank is the layer that makes the design usable from software.} It catches
the pulses and holds them as levels that you have to clear yourself. Every "write \code{0x1} here
to acknowledge" in the register map is that mechanism.
